% !PRINTED 549
\weberpage{549}
Wir wollen eine positive Grösse $\varrho_1$ und einen Winkel $\Theta_1$ so
annehmen, dass
\weq{9}{\varphi_1 = \varrho_1 e^{i\Theta_1}, \qquad \varphi_1' = \varrho_1 e^{-i\Theta_1},}
oder
\weq{10}{\begin{aligned} \varphi_1 &= \varrho_1(\cos\Theta_1 + i\sin\Theta_1) \\ \varphi_1' &= \varrho_1(\cos\Theta_1 - i\sin\Theta_1), \end{aligned}}
woraus noch folgt:
\weq{11}{\varrho_1^2 = \varphi_1 \varphi_1'.}
Nun ist
\weq{12}{(\varepsilon^{\lambda_1}, \alpha)(\varepsilon^{-\lambda_1}, \alpha) = \pm a_1}
eine Grösse des Körpers $\Omega(\varepsilon)$ nach (\S.\,164, 3.), und zwar ist es,
da sie sich beim Uebergang zum conjugirt imaginären Werth,
d. h. bei der Vertauschung von $\varepsilon$ und $\varepsilon^{-1}$ nicht ändert, eine
reelle Grösse. Das Vorzeichen wollen wir so bestimmen, dass $a_1$
positiv ist. Es ist also nach (11), da $\varrho_1^2$ positiv ist,
\[
\varrho_1^2 = (\pm a_1)^m = a_1^m,
\]
und es ergiebt sich daraus, dass bei ungeradem $m$ jedenfalls das
obere Zeichen gilt; bei geradem $m$ kann auch das untere eintreten. Es ist dann
\weq{13}{\varrho_1 = \sqrt{a_1^m},}
wo die Quadratwurzel positiv zu nehmen ist.

Ferner sind aus gleichem Grunde
\weq{14}{\frac{\varphi_1 + \varphi_1'}{2} = b_1, \qquad \frac{\varphi_1 - \varphi_1'}{2i} = c_1}
reelle Grössen in $\Omega(\varepsilon)$, und es ergiebt sich aus (10)
\weq{15}{\cos\Theta_1 = \frac{b_1}{\sqrt{a_1^m}}, \qquad \sin\Theta_1 = \frac{c_1}{\sqrt{a_1^m}},}
woraus noch die Relation folgt:
\[
a_1^m = b_1^2 + c_1^2.
\]

Demnach ergiebt sich
\[
\sqrt[p_1]{\varphi_1} = \sqrt{a_1^{m_1}}\, e^{\frac{i\Theta_1}{p_1}};
\]
und nun verfährt man mit den Functionen $\varphi_2$ u. s. f. ebenso.
Man bestimmt also $\Theta_2 \ldots$ aus einem System von Gleichungen
wie (15) und erhält, wenn man noch
\weq{16}{m_1\Theta_1 + m_2\Theta_2 + \cdots = \Theta}
setzt, nach (8)
%CONT>>
